\chapter{A Small System with Real Boundaries}

Rails applications are easy to test when every dependency lives in the same
process or on the same development machine. Production systems are less tidy.
They cross network boundaries, wait for other services to become ready, and
must behave sensibly when a dependency is unavailable.

In this tutorial we will build a small multi-service Rails system on FreeBSD.
Each service will run in a BSD jail managed through BSDRunner and Blackship.
The finished system will let us exercise a complete request across real
operating-system and network boundaries.

\section{The system we will build}

Our initial topology has four service jails:

\begin{description}
  \item[web] An HTTP reverse proxy and the only service intended for inbound
    access from the host.
  \item[store] The primary Rails application. It renders the user-facing
    workflow and calls the catalog API over HTTP.
  \item[catalog] A Rails API that owns product and availability data.
  \item[database] PostgreSQL, with separate logical databases and credentials
    for the two Rails applications.
\end{description}

The first end-to-end test will enter through the reverse proxy, visit a page in
the store, and verify that data obtained from the catalog API appears in the
response. Later, we will stop the catalog jail deliberately and test the
store's failure behavior.

\begin{lstlisting}
test -> web -> store -> catalog
                 |         |
                 +----+----+
                      |
                  database
\end{lstlisting}

The diagram is deliberately a destination, not our first task. We will add one
boundary at a time and keep the system runnable after every chapter.

\section{What this tutorial emphasizes}

The goal is not merely to make Rails boot inside a jail. We will use the
project to learn how a small distributed system is assembled and tested:

\begin{itemize}
  \item repeatable jail creation and service configuration;
  \item private addressing and explicit exposure through the proxy;
  \item database ownership, migrations, and credentials;
  \item dependency ordering and readiness rather than arbitrary delays;
  \item end-to-end tests that cross service boundaries; and
  \item observable, intentional failure when a dependency disappears.
\end{itemize}

Unless a chapter says otherwise, commands are run on the FreeBSD host. Actions
that change the host or a jail will be identified before they are performed.
The examples favor small, inspectable configuration over hidden automation.

\section{Exercise 1: Build this handbook}

Before changing the jail environment, verify that the tutorial source can
produce a PDF. From the BSDRunner repository, run:

\begin{lstlisting}[language=sh]
cd docs/rails-on-jails
make
\end{lstlisting}

The command should finish successfully and create \texttt{build/main.pdf}.
Open the file and confirm that this chapter and its topology diagram are
present.

\subsection*{Checkpoint}

Do not provision a jail yet. At the end of this exercise you should have only:

\begin{itemize}
  \item a successfully generated PDF;
  \item a shared picture of the four-service topology; and
  \item no changes to the running FreeBSD jail environment.
\end{itemize}

The next exercise will inventory the host, Blackship, ZFS storage, and the
BSDRunner private network. That evidence will determine the exact provisioning
steps instead of assuming that every FreeBSD host begins in the same state.
